\documentclass[11pt,a4paper]{article}
\usepackage[margin=27mm]{geometry}
\usepackage{amsmath,amssymb,amsthm,mathtools,booktabs,microtype}
\usepackage[colorlinks=true,allcolors=blue]{hyperref}
\setlength{\emergencystretch}{2em}
\newtheorem{theorem}{Theorem}
\newtheorem{proposition}[theorem]{Proposition}
\newtheorem{lemma}[theorem]{Lemma}
\newtheorem{remark}[theorem]{Remark}
\newtheorem{corollary}[theorem]{Corollary}
\newcommand{\R}{\mathbb R}
\newcommand{\C}{\mathbb C}
\newcommand{\conv}{\operatorname{conv}}
\newcommand{\Tr}{\operatorname{Tr}}
\newcommand{\supp}{\operatorname{supp}}
\newcommand{\one}{\mathbf 1}
\newcommand{\Hinf}{H_{\infty}}
\newcommand{\Finf}{F_{\theta,\infty}}
\newcommand{\Ginf}{G_{\xi,\infty}}
\newcommand{\ket}[1]{|#1\rangle}
\title{A four-qubit criterion for entropy minimax duality\\
\large Minimality, sharp R\'enyi threshold and a W-state counterexample}
\author{Lluis Eriksson\\ \small Independent researcher}
\date{17 October 2026\\ \small Research working paper}
\begin{document}
\maketitle
\begin{abstract}
We characterize entropy functions for which support--radial minimax
duality holds on every multiqubit orbit closure. For a continuous,
strictly increasing function of the smaller marginal eigenvalue,
vanishing at zero, universal validity is equivalent to concavity.
It suffices to test a single tensor, the four-qubit W state, over all
positive weights. Four qubits are minimal even within the R\'enyi family:
order three satisfies the identity for every tensor on at most
three qubits, but fails for the four-qubit W state. Order four
already fails with three qubits. For binary R\'enyi entropy this gives the sharp
range \(0<\alpha\le2\); every order above two admits a four-qubit
obstruction. The scalar curvature facts are established properties;
the proposed contribution is their universal W-state test.
At min-entropy, the three-qubit W tensor already contradicts the first
identity of Theorem 60 in Christandl--Lysikov--Zuiddam,
both in arXiv:2012.14412v3 and in the 2023 journal version:
its two sides are \(3/2\) and \(2^{2/3}\).
We compute the weighted min-entropy functionals of all W tensors,
repair the identity by convexifying entropy vectors, and retain the
reverse identity. A typical-projector argument explains tensor
regularization in terms of established Shannon functionals; the
unweighted identification was already known from G-stable rank.
The proofs include exact finite-order witnesses. Priority and
human validation remain open.
\end{abstract}

\section{Definitions and the exact scope of the obstruction}
Let \(T\ne0\) be a complex tensor in \(\bigotimes_{i=1}^k\C^{d_i}\).
Write \(\Pi(T)\) for the tuples of ordered reduced-density-matrix
spectra of normalized representatives of its projective local
general-linear orbit closure. This is a compact set. Logarithms are
base two. For a probability vector \(p\), let
\[
 \Hinf(p)=-\log\max_jp_j,\qquad H(p)=-\sum_jp_j\log p_j .
\]
Set
\[
 \begin{gathered}
 \Theta_k=\{\theta\in\R_{\ge0}^k:\textstyle\sum_i\theta_i=1\},\qquad
 \Xi_k=\{\xi\in\R_{\ge0}^k:\max_i\xi_i=1\},\\
 \Theta_k(\xi)=\{\theta\in\Theta_k:\supp\theta\subseteq\supp\xi\}.
 \end{gathered}
\]
The natural \(k\)-party versions of the min-entropy quantities are
\begin{align}
 F_{\theta,\infty}(T)
 &=\max_{p\in\Pi(T)}2^{\sum_i\theta_i\Hinf(p_i)},\label{eq:F}\\
 G_{\xi,\infty}(T)
 &=\max_{p\in\Pi(T)}\min_{i:\xi_i>0}2^{\Hinf(p_i)/\xi_i}.
 \label{eq:G}
\end{align}
For \(k=3\), these are the definitions in \cite[Section 5.3]{CLZ}.
The first identity of its Theorem 60, also stated in its introductory
Theorem 14, asserts
\begin{equation}
 G_{\xi,\infty}(T)
 \stackrel{\text{claimed}}{=}
 \min_{\theta\in\Theta_3(\xi)}
 F_{\theta,\infty}(T)^{1/\langle\theta,\xi\rangle}.
 \label{eq:claimed}
\end{equation}
We assessed \textbf{arXiv:2012.14412v3, dated 19 April 2023}, with
Theorem 60 on printed page 26, and the \textbf{final journal version},
J. Math. Pures Appl.\ 172 (2023), with that theorem on page 327.
The journal text was checked in the University of Amsterdam repository
copy linked in \cite{CLZ}; it contains the same first identity and
the same preceding concavity assertion. This comparison does not
establish the absence of a later correction elsewhere.

The general concave minimax theorem in \cite[Theorem 50]{CLZ} is not
contradicted. The issue is its use for min-entropy: \(\Hinf\) is not
concave. For \(p=(1,0)\), \(q=(1/2,1/2)\),
\[
 \Hinf((p+q)/2)=\log(4/3)<\tfrac12
 =\tfrac12\Hinf(p)+\tfrac12\Hinf(q).
\]
Quasiconcavity does not give the concavity required for that application.
The established Shannon-entropy identities are unaffected.

\section{Exact W-state formulas}
Let
\[
 W_k=\sum_{j=1}^k\ket{0\cdots 0\,1_j\,0\cdots0},\qquad k\ge2,
\]
and write \(\theta_{(1)}\ge\theta_{(2)}\ge\cdots\) for sorted weights.

\begin{lemma}\label{lem:orbit}
For every marginal tuple in \(\Pi(W_k)\), its largest eigenvalues
\(\lambda_i\) satisfy
\begin{equation}
 \tfrac12\le\lambda_i\le1,\qquad \sum_i\lambda_i\ge k-1.
 \label{eq:lambda}
\end{equation}
Conversely, every tuple with \(\lambda_i\in[1/2,1]\) and
\(\sum_i\lambda_i=k-1\) is attained in the orbit closure.
\end{lemma}
\begin{proof}
For a tensor in the orbit, apply local unitaries that take each
\(g_i\ket0\) to a multiple of \(\ket0\). Expanding each \(g_i\ket1\)
then gives the normalized form
\[
 \psi=a_0\ket{0\cdots0}+\sum_{j=1}^ka_j\ket{0\cdots1_j\cdots0},
 \qquad \sum_{j=0}^k|a_j|^2=1.
\]
With \(r_j=|a_j|^2\), the \(j\)-th marginal has diagonal entries
\(1-r_j,r_j\). Its smaller eigenvalue \(\mu_j\) is at most \(r_j\),
by the Rayleigh principle. Hence
\(\sum_j\mu_j\le\sum_jr_j\le1\), giving \eqref{eq:lambda};
continuity extends this bound to the closure.

For the converse, put \(r_j=1-\lambda_j\), so that
\(0\le r_j\le1/2\) and \(\sum_jr_j=1\). The state
\(\sum_j\sqrt{r_j}\ket{0\cdots1_j\cdots0}\) has diagonal marginals
\((1-r_j,r_j)\). Positive coefficients are obtained by invertible
diagonal local maps on \(W_k\), and zero coefficients by limits.
\end{proof}

\begin{theorem}\label{thm:W}
For every \(\theta\in\Theta_k\),
\begin{equation}
 F_{\theta,\infty}(W_k)=2^{\theta_{(1)}+\theta_{(2)}}.
 \label{eq:Wf}
\end{equation}
For \(\xi\in\Xi_k\), write \(I=\supp\xi\) and \(s=|I|\).
If \(s=1\), then \(G_{\xi,\infty}(W_k)=2\).
If \(s\ge2\), let \(t_*>1\) be the unique solution of
\[
 \sum_{i\in I}t_*^{-\xi_i}=s-1.
\]
Then
\begin{equation}
 G_{\xi,\infty}(W_k)=\min\{2,t_*\}.
 \label{eq:Wg}
\end{equation}
In particular,
\begin{equation}
 G_{\one,\infty}(W_k)=\frac{k}{k-1},\qquad
 \min_{\theta\in\Theta_k}F_{\theta,\infty}(W_k)=2^{2/k}.
 \label{eq:unweighted}
\end{equation}
\end{theorem}
\begin{proof}
Put \(\mu_i=1-\lambda_i\). Lemma \ref{lem:orbit} places this tuple
in \(Q=\{\mu:0\le\mu_i\le1/2,\ \sum_i\mu_i\le1\}\).
The function \(-\sum_i\theta_i\log(1-\mu_i)\) is convex on \(Q\).
Its maximum is attained at a vertex. The vertices of \(Q\) are
the zero vector, the vectors with one entry \(1/2\), and those with
two entries \(1/2\). Indeed, after scaling by two, a fractional
vertex would admit a coordinate perturbation, or a two-coordinate
perturbation when the sum constraint is tight.
The largest objective value is \(\theta_{(1)}+\theta_{(2)}\).
It is attained by projecting all but the corresponding two parties
onto \(\ket0\), leaving an entangled pair and pure spectators.
This proves \eqref{eq:Wf}.

For \(G_{\xi,\infty}\ge t\), active largest marginal eigenvalues
must satisfy \(\lambda_i\le t^{-\xi_i}\); for inactive coordinates
this is the automatic bound \(\lambda_i\le1=t^0\). Since some \(\xi_i=1\),
the necessary conditions are
\begin{equation}
 1\le t\le2,\qquad \sum_it^{-\xi_i}\ge k-1.
 \label{eq:feasible}
\end{equation}
They are sufficient: choose
\(\lambda_i\in[1/2,t^{-\xi_i}]\) with sum \(k-1\).
Such a choice exists because the lower endpoints sum to \(k/2\le k-1\),
whereas the upper endpoints sum to at least \(k-1\).
Use the converse in Lemma \ref{lem:orbit}. The sum in
\eqref{eq:feasible} is equivalently
\(\sum_{i\in I}t^{-\xi_i}\ge s-1\).
For \(s=1\) it imposes no additional restriction and the value is two.
For \(s\ge2\), the active sum decreases strictly and continuously
from \(s\) to zero, proving \eqref{eq:Wg}. In particular \(s=2\)
also gives value two, since each \(2^{-\xi_i}\ge1/2\).

The unweighted case follows by symmetry. The sum of the two largest
entries of a probability vector is at least \(2/k\), with equality
at the uniform vector, proving the second formula.
\end{proof}

\begin{proposition}[Counterexample to \eqref{eq:claimed}]\label{prop:counter}
For \(W=W_3=\ket{100}+\ket{010}+\ket{001}\) and \(\xi=(1,1,1)\),
\[
 G_{\xi,\infty}(W)=\frac32
 <2^{2/3}
 =\min_{\theta\in\Theta_3}F_{\theta,\infty}(W)^{1/\langle\theta,\xi\rangle}.
\]
\end{proposition}
\begin{proof}
Here \(\langle\theta,\xi\rangle=1\). Apply Theorem \ref{thm:W};
the strict inequality is exactly \(27/8<4\) after cubing.
\end{proof}

The value \(3/2\) is already known from Derksen's computation of the
G-stable rank \cite{Derksen}; the W-polytope geometry is likewise
established \cite{Walter}. We use a self-contained orbit argument to
make the contradiction independently checkable. For every \(k>2\),
the analogous gap is also strict: \((1-1/k)^k>1/4\), since
\(x\log(1-1/x)\) increases for \(x>1\) and has value \(\log(1/4)\)
at \(x=2\).

\section{What the minimax expression actually computes}
The distinction is between an entropy image and its convex hull.
The following statement does not assume that the original parameter
space is convex.

\begin{theorem}[Convexified radial duality]\label{thm:convex}
Let \(E\subset\R_{\ge0}^k\) be nonempty and compact. Define
\[
 f_\theta=\max_{e\in E}\langle\theta,e\rangle,\quad
 g_\xi=\max_{e\in E}\min_{i:\xi_i>0}\frac{e_i}{\xi_i},\quad
 \widehat g_\xi=\max_{z\in\conv E}
                    \min_{i:\xi_i>0}\frac{z_i}{\xi_i}.
\]
Then
\begin{equation}
 \widehat g_\xi
 =\min_{\theta\in\Theta_k(\xi)}
       \frac{f_\theta}{\langle\theta,\xi\rangle},
 \qquad g_\xi\le\widehat g_\xi.
 \label{eq:repair}
\end{equation}
An optimizer for the active coordinates of \(\widehat g_\xi\)
can be obtained from a convex combination of at most
\(|\supp\xi|\) points of \(E\). Moreover, without convexifying,
\begin{equation}
 f_\theta=\max_{\xi\in\Xi_k}g_\xi\langle\theta,\xi\rangle.
 \label{eq:reverse}
\end{equation}
\end{theorem}
\begin{proof}
Let \(I=\supp\xi\) and
\(\Lambda=\{u\ge0:\supp u\subseteq I,\ \langle u,\xi\rangle=1\}\).
This is compact and convex, and
\(\min_{u\in\Lambda}\langle u,z\rangle=\min_{i\in I}z_i/\xi_i\).
The bilinear minimax theorem gives
\[
 \max_{z\in\conv E}\min_{u\in\Lambda}\langle u,z\rangle
 =\min_{u\in\Lambda}\max_{e\in E}\langle u,e\rangle.
\]
Parameterize \(u=\theta/\langle\theta,\xi\rangle\) to obtain
\eqref{eq:repair}. This is an application of classical convex
minimax, not a new minimax principle.

For the support bound, project onto the coordinates in \(I\).
Let \(z_*\) and \(u_*\) be primal and dual optimizers, with common
value \(v\). The saddle-point equalities follow from
\[
 v=\min_{u\in\Lambda}\langle u,z_*\rangle
 \le\langle u_*,z_*\rangle
 \le\max_{z\in\conv E}\langle u_*,z\rangle=v.
\]
Thus \(z_*\) lies in the supporting face exposed by \(u_*\).
That face has affine dimension at most \(|I|-1\), since
\(u_*\ne0\). Carath\'eodory's theorem on that face gives at most
\(|I|\) points. Those points may be taken from the projected \(E\):
a convex combination can attain the supporting value only through
points attaining that value.

For \eqref{eq:reverse}, an optimizer \(e\) of \(g_\xi\) satisfies
\(e_i\ge g_\xi\xi_i\) for all \(i\), including the inactive coordinates.
Thus \(g_\xi\langle\theta,\xi\rangle\le f_\theta\).
Conversely, choose \(e\) attaining \(f_\theta\). If \(e\ne0\), set
\(a=\max_i e_i\) and \(\xi=e/a\). Then \(g_\xi\ge a\), giving the
opposite inequality. If \(f_\theta=0\), nonnegativity and the first
inequality already prove equality.
\end{proof}

Apply Theorem \ref{thm:convex} to
\[
 E_\infty(T)=\{(\Hinf(p_1),\ldots,\Hinf(p_k)):p\in\Pi(T)\}.
\]
The right-hand side of \eqref{eq:claimed} is \(2^{\widehat g_\xi}\),
whereas its left-hand side is \(2^{g_\xi}\). Thus convexification
repairs precisely the failed identity. The second identity in
\cite[Theorem 60]{CLZ}, namely the multiplicative form of
\eqref{eq:reverse}, remains valid.

For \(W_k\), the downward closure of \(\conv E_\infty(W_k)\) in
the nonnegative orthant is
\[
 \{z\in[0,1]^k:\textstyle\sum_i z_i\le2\}.
\]
Indeed, the upper bounds follow from \eqref{eq:Wf}, and the pair
indicator vectors are attained by entangled pairs with pure spectators.
Their convex hull is the face \(\sum_i z_i=2\); its downward closure
is the displayed set. Consequently,
\begin{equation}
 \widehat g_\xi=\min\{1,\,2/\textstyle\sum_i\xi_i\},
 \qquad \xi\in\Xi_k.
 \label{eq:Wconvex}
\end{equation}
For \(W_3\), averaging the three entropy vectors
\((0,1,1),(1,0,1),(1,1,0)\) gives the balanced vector
\((2/3,2/3,2/3)\) responsible for the incorrect one-copy prediction.

Other basic properties do not depend on the failed identity.
Orbit-closure monotonicity is immediate from inclusion; on the
diagonal tensor of size \(r\), \(G_{\xi,\infty}=r\).
For tensor products, local maximum eigenvalues multiply. Taking
optimizing states and using
\(\min_i(a_i+b_i)\ge\min_i a_i+\min_i b_i\) gives
\[
 G_{\xi,\infty}(S\otimes T)
 \ge G_{\xi,\infty}(S)\,G_{\xi,\infty}(T).
\]
Thus the monotonicity, normalization and supermultiplicativity
assertions of \cite[Corollary 61]{CLZ} have direct proofs.

\section{A universal four-qubit test}\label{sec:criterion}
The obstruction extends beyond min-entropy. Let
\(h:[0,1/2]\to\R_{\ge0}\) be continuous and strictly increasing,
with \(h(0)=0\). For a nonzero multiqubit tensor \(T\), let
\(P(T)\) denote its orbit-closure marginal tuples in smaller-eigenvalue
coordinates \(\mu=(\mu_1,\ldots,\mu_k)\). This is an affine image of
the entanglement polytope and hence is compact and convex
\cite[Supplementary Information, Corollary 2]{Walter}.
Define logarithmic quantities
\[
 f^h_\theta(T)=\max_{\mu\in P(T)}\sum_i\theta_i h(\mu_i),\qquad
 g^h_\xi(T)=\max_{\mu\in P(T)}\min_{i:\xi_i>0}\frac{h(\mu_i)}{\xi_i}.
\]

\begin{theorem}[Four-qubit criterion]\label{thm:criterion}
Under the stated assumptions on \(h\), the following are equivalent:
\begin{enumerate}
 \item \(h\) is concave on \([0,1/2]\).
 \item For every \(k\ge2\), every nonzero \(k\)-qubit tensor \(T\),
 and every \(\xi\in\Xi_k\),
 \begin{equation}
  g^h_\xi(T)=\min_{\theta\in\Theta_k(\xi)}
       \frac{f^h_\theta(T)}{\langle\theta,\xi\rangle}.
  \label{eq:hduality}
 \end{equation}
 \item Equation \eqref{eq:hduality} holds for \(T=W_4\) and every
 \(\xi\in\Xi_4\) with strictly positive coordinates.
\end{enumerate}
If concavity fails, a witness to failure on \(W_4\) can be obtained
by mixing just two entropy vectors of weighted W states.
\end{theorem}
\begin{proof}
Suppose \(h\) is concave. On the compact convex set \(P(T)\),
the function \(\sum_i u_i h(\mu_i)\) is concave in \(\mu\) for
every \(u\ge0\), and is linear in \(u\). Apply the concave minimax
theorem with the compact weight set \(\Lambda\) from the proof of
Theorem \ref{thm:convex}. This proves \eqref{eq:hduality}, including
weights with inactive coordinates. The second assertion plainly
implies the third.

For the converse, continuity implies that failure of concavity
has a midpoint witness: there are \(0\le x<y\le1/2\) such that,
writing \(a=(x+y)/2\) and \(b=1/2-a\),
\[
 c=\frac{h(x)+h(y)}2>h(a).
\]
Both \(a\) and \(b\) lie strictly between zero and \(1/2\).
Lemma \ref{lem:orbit} shows that
\(r^+=(x,y,b,b)\) and \(r^-=(y,x,b,b)\) belong to \(P(W_4)\):
their coordinates lie in \([0,1/2]\) and sum to one.
Their entropy vectors have average
\(z=(c,c,h(b),h(b))\). Set \(d=\max_i z_i>0\) and \(\xi=z/d\).
These weights are strictly positive. The convexified radial value
is at least \(d\), since \(z\) lies in the convex hull.

The original radial value is strictly less than \(d\).
Otherwise, compactness would give \(\mu\in P(W_4)\) with
\(h(\mu_i)\ge z_i\) for every \(i\). Strict monotonicity would imply
\(\mu_1>a\), \(\mu_2>a\), \(\mu_3\ge b\), \(\mu_4\ge b\),
and thus \(\sum_i\mu_i>2a+2b=1\), contradicting
Lemma \ref{lem:orbit}. By Theorem \ref{thm:convex}, the right-hand
side of \eqref{eq:hduality} equals the convexified radial value,
so \eqref{eq:hduality} fails for these weights.
\end{proof}

The theorem concerns universal validity over all weights and all
multiqubit orbit closures. It does not say that every choice of
weights fails when \(h\) is nonconcave. The sufficiency direction
is the usual concave minimax argument. The converse identifies
\(W_4\) as a single test tensor for this whole class of scalar functions.

\subsection{Minimality within the R\'enyi family}
Four qubits are necessary even if the test functions are restricted to
binary R\'enyi entropies. This is stronger than minimality witnessed by
an arbitrary increasing scalar function.

\begin{theorem}[Minimal number of qubits]\label{thm:minimality}
The binary order-three R\'enyi entropy
\[
 h_3(p)=-\frac12\log(1-3p+3p^2),\qquad 0\le p\le\frac12,
\]
satisfies \eqref{eq:hduality} for every nonzero tensor on two or three
qubits and every admissible weight vector. It fails on \(W_4\) with
the positive weights specified in Section~\ref{sec:order-three}.
Consequently four is the minimal universal testing size in
Theorem~\ref{thm:criterion}, including within the R\'enyi family.
\end{theorem}
\begin{proof}
Write \(h=h_3\), \(\phi=h^{-1}\). The function \(h\) is continuous
and strictly increasing, with \(h(0)=0\) and \(h(1/2)=1\).
By Lemma~\ref{lem:orbit}, the downward closure of its entropy image
on \(W_3\) is exactly
\[
 D=\{z\in[0,1]^3:\phi(z_1)+\phi(z_2)+\phi(z_3)\le1\}.
\]
For sufficiency, increase any inverse-coordinate tuple of sum at most
one, keeping each coordinate at most \(1/2\), until its sum is one.
The remaining total capacity is \(3/2-\sum_i p_i\ge1-\sum_i p_i\).
The resulting tuple is realized by a weighted W state and dominates
the original entropy vector. Necessity is the marginal sum bound.

We show that \(D\) is convex, although \(h\) is not concave. For
\(0<p_i<1/2\), \(\sum_i p_i=1\), set \(x_i=1-2p_i>0\), so
\(\sum_i x_i=1\). Direct differentiation gives
\[
 A(p_i):=-\frac{h''(p_i)}{h'(p_i)}
 =\frac{2(1-3x_i^2)}{x_i(1+3x_i^2)}.
\]
For every real tuple \(v\) with \(\sum_i v_i=0\), we claim
\begin{equation}
 \sum_{i=1}^3 A(p_i)v_i^2\ge0.\label{eq:minimal-curvature}
\end{equation}
If all \(x_i\le1/\sqrt3\), each coefficient is nonnegative.
Otherwise exactly one, say \(x=x_1\), exceeds \(1/\sqrt3\), so
\(A(p_1)<0<A(p_2),A(p_3)\). Define
\[
 B(x)=\frac{x(1+3x^2)}{1-3x^2}.
\]
On \([0,1/\sqrt3)\), the ratio \(B(x)/x\), extended at zero,
is increasing: its derivative is \(12x/(1-3x^2)^2\).
Thus \(B(x_2)+B(x_3)\le B(x_2+x_3)=B(1-x)\).
With \(s=x(1-x)\le1/4\), expansion yields
\[
 B(x)+B(1-x)
 =\frac{4-12s-9s^2}{(1-3x^2)(1-3(1-x)^2)}<0.
\]
The numerator is at least \(7/16\); the denominator is negative.
Since \(B(x_i)=2/A(p_i)\), it follows that
\[
 A(p_1)+\frac{A(p_2)A(p_3)}{A(p_2)+A(p_3)}>0.
\]
For fixed \(v_1\) and \(v_2+v_3=-v_1\), the minimum of
\(A(p_2)v_2^2+A(p_3)v_3^2\) is the second coefficient in this
display times \(v_1^2\). This proves \eqref{eq:minimal-curvature}.

For global convexity, write \(D\) as \(0\le z_3\le g(z_1,z_2)\), where
\[
 g(z_1,z_2)=
 \begin{cases}
 1,&\phi(z_1)+\phi(z_2)\le1/2,\\
 h(1-\phi(z_1)-\phi(z_2)),&\phi(z_1)+\phi(z_2)>1/2.
 \end{cases}
\]
First take a line segment in the open square \((0,1)^2\), where
\(\phi\) is analytic with finite derivatives. On curved pieces put
\(p_i=\phi(z_i)\), \(i=1,2\), and \(p_3=1-p_1-p_2\).
Then \(p_i''=A(p_i)(p_i')^2\) for \(i=1,2\), so
\[
 g''=-h'(p_3)\sum_{i=1}^3 A(p_i)(p_i')^2\le0.
\]
Flat pieces have zero curvature. At the interface \(p_3=1/2\),
both one-sided derivatives of \(g\) are zero because \(h'(1/2)=0\).
Analyticity gives finitely many interface contacts on a compact
segment unless its interface equation is identically satisfied;
in that case \(g\) is constant. Each line restriction is therefore
concave. Continuity extends concavity to the closed square, without
requiring a finite derivative of \(\phi\) at its upper endpoint.
Hence \(D\) is convex. Nonnegative support and radial maxima do not
change under downward closure, and Theorem~\ref{thm:convex} proves
the identity for the W class.

The established pure-state SLOCC classification \cite{DVC} covers
all remaining three-qubit tensors: every nonzero tensor can be
normalized as a pure state, and local invertible equivalence
preserves its projective orbit closure. The GHZ class contains a state
with all smaller marginal eigenvalues equal to \(1/2\). Biseparable
classes contain a Bell pair and a pure spectator, and the product
class gives a point. Their entropy downsets are a box, a box face
or a point, so they are convex. For two qubits the product and
Schmidt-rank-two classes likewise give a point or a box.
The explicit four-qubit failure is proved in
Section~\ref{sec:order-three} by an exact rational midpoint gap.
\end{proof}

This is minimality for universal testing, not a statement that every
nonconcave entropy needs four qubits. Indeed, order four already fails
on \(W_3\). To see this, take the weighted W states with marginal
coordinates \((1/5,3/10,1/2)\) and \((3/10,1/5,1/2)\).
For \(f_4(p)=p^4+(1-p)^4\), the rational identity
\[
 f_4(1/4)^2-f_4(1/5)f_4(3/10)=\frac{3457649}{6400000000}>0
\]
implies \(c=(h_4(1/5)+h_4(3/10))/2>h_4(1/4)\).
Their entropy average is \((c,c,1)\). No original W-state entropy
vector dominates it: its third inverse coordinate would be at least
\(1/2\), and its first two would each exceed \(1/4\).
Thus \(\xi=(c,c,1)\) gives a strict gap. We do not classify the
complete three-qubit transition as the R\'enyi order varies.

\subsection{The sharp R\'enyi range}
Write
\[
 h_\alpha(p)=\frac{\log\bigl(p^\alpha+(1-p)^\alpha\bigr)}{1-\alpha},
 \quad 0<\alpha<\infty,\quad\alpha\ne1,
\]
with the continuous Shannon limit at \(\alpha=1\), and
\(h_\infty(p)=-\log(1-p)\) for \(0\le p\le1/2\).
These functions are continuous, strictly increasing on this interval,
and vanish at zero. We include the scalar curvature argument so that
the threshold does not depend on a numerical plot. The loss of
concavity for orders above two is already recorded, for example,
in \cite[Lemma 4.3]{Hillion}; that scalar observation is not new.

\begin{lemma}[Binary R\'enyi curvature]\label{lem:renyi}
The function \(h_\alpha\) is concave on \([0,1/2]\) exactly when
\(0<\alpha\le2\). For every finite \(\alpha>2\),
\[
 \lim_{p\downarrow0}h_\alpha''(p)
 =\frac{\alpha}{(\alpha-1)\ln2}>0.
\]
The min-entropy limit is strictly convex on this interval.
\end{lemma}
\begin{proof}
For \(0<\alpha<1\), the sum \(p^\alpha+(1-p)^\alpha\) is concave
and positive. Composition with the increasing concave logarithm,
and division by \(1-\alpha>0\), preserve concavity. Shannon's
binary entropy is concave as well.

For \(1<\alpha\le2\), set \(F=p^\alpha+(1-p)^\alpha\),
\(r=\alpha-1\in(0,1]\), and \(t=p/(1-p)>0\). Expansion gives
\[
 \frac{FF''-(F')^2}{\alpha[p(1-p)]^{\alpha-1}}
 =r(t+t^{-1})+2\alpha-(t^r+t^{-r}).
\]
Convexity of \(s\mapsto t^s+t^{-s}\) on \([0,1]\) gives
\(t^r+t^{-r}\le r(t+t^{-1})+2(1-r)\).
Thus the displayed quantity is at least \(4r>0\), and
\[
 h_\alpha''(p)
 =-\frac{FF''-(F')^2}{(\alpha-1)(\ln2)F^2}<0.
\]
Continuity extends concavity to the endpoints. For \(\alpha>2\),
as \(p\downarrow0\), \(F\to1\), \(F'\to-\alpha\), and
\(F''\to\alpha(\alpha-1)\), giving the asserted positive limit.
Finally, \(h_\infty''(p)=1/((1-p)^2\ln2)>0\).
\end{proof}

\begin{corollary}[Sharp multiqubit threshold]\label{cor:threshold}
The support--radial minimax identity formed from R\'enyi marginal
entropies is valid on every multiqubit orbit closure, for every weight
vector, if and only if \(0<\alpha\le2\).
For every \(2<\alpha\le\infty\), it already fails for \(W_4\)
with some strictly positive weights.
\end{corollary}
\begin{proof}
Apply Theorem \ref{thm:criterion} and Lemma \ref{lem:renyi}.
\end{proof}

This is a statement about the marginal support and radial quantities
defined here. It is not a statement about conditional R\'enyi entropy,
LOCC spectral functionals, or general entropic inequalities.
In particular, no claim about the Shepp--Olkin conjectures for sums
of Bernoulli variables is implied.

\subsection{An exact witness at order three}\label{sec:order-three}
Take \(h=h_3\), \(x=1/20\), \(y=3/20\), \(a=1/10\),
and \(b=2/5\). Since \(h_3(p)=-\tfrac12\log f(p)\) for
\(f(p)=1-3p+3p^2\), the required strict midpoint gap follows from
the rational identity
\[
 f(a)^2-f(x)f(y)=\frac{543}{160000}>0,
 \qquad f(x)f(y)=\frac{84721}{160000}.
\]
The two normalized states have squared single-excitation amplitudes
\((1,3,8,8)/20\) and \((3,1,8,8)/20\), so their reduced spectra
and membership in the W orbit are explicit. With
\(c=(h_3(x)+h_3(y))/2\), choose
\[
 \xi=(c/h_3(b),c/h_3(b),1,1).
\]
Theorem \ref{thm:criterion} gives a strict duality gap with these
exactly specified weights.

For numerical orientation, the first weight is approximately
\(0.2497358423\). The original logarithmic radial value is
approximately \(0.9169557149\), whereas the convexified value is
at least \(h_3(2/5)\approx0.9182506339\).
The first value can be computed from
\[
 \sum_{i=1}^4 h_3^{-1}(t\xi_i)=1,
 \qquad h_3^{-1}(u)=\frac{3-\sqrt{12\,2^{-2u}-3}}6.
\]
Indeed, feasibility requires \(\mu_i\ge h_3^{-1}(t\xi_i)\) and
\(\sum_i\mu_i\le1\); it is sufficient by filling the coordinates
to sum to one, each at most \(1/2\), and using Lemma \ref{lem:orbit}.
The decimal values are high-precision diagnostics, not interval
certificates. Strictness of the gap rests on the rational identity
and the theorem's proof.

\section{Tensor regularization and simultaneous spectral flattening}
Here we restrict the stated identification to tripartite tensors,
so that all external inputs agree directly with \cite{CLZ,CVZ}.
Let
\[
 F_\theta(T)=\max_{p\in\Pi(T)}2^{\sum_i\theta_iH(p_i)},\qquad
 G_\xi^{H}(T)=\max_{p\in\Pi(T)}\min_{i:\xi_i>0}2^{H(p_i)/\xi_i}.
\]
We use two established facts: the Shannon quantum functional
\(F_\theta\) is multiplicative \cite[Theorem 3.30 and Corollary 3.31]{CVZ}, and the Shannon minimax
identity is valid \cite{CLZ}. These are inputs, not conclusions
claimed as new in this note.

\begin{lemma}[Simultaneous local typical projection]\label{lem:typical}
For a normalized multipartite pure state \(\psi\), let \(H_i\)
be the entropy of its \(i\)-th marginal. For every \(\delta>0\)
there are local projectors on \(\psi^{\otimes m}\), with success
probability \(q_m\to1\), whose normalized output \(\sigma^{(m)}\)
satisfies, for every party,
\[
 \Hinf(\sigma_i^{(m)})\ge m(H_i-\delta)+\log q_m.
 \label{eq:flatten}
\]
\end{lemma}
\begin{proof}
Let \(P_i^{(m)}\) be the spectral projector of \(\rho_i^{\otimes m}\)
onto eigenvalues at most \(2^{-m(H_i-\delta)}\).
The law of large numbers for the random variable \(-\log p\), sampled
according to the positive eigenvalues of \(\rho_i\), gives
\(\Tr(P_i^{(m)}\rho_i^{\otimes m})\to1\).
The local projectors commute, so for \(P=\bigotimes_iP_i^{(m)}\),
\[
 q_m=\|P\psi^{\otimes m}\|^2
 \ge1-\sum_i\bigl(1-\Tr(P_i^{(m)}\rho_i^{\otimes m})\bigr)\to1.
\]
For sufficiently large \(m\), normalize \(P\psi^{\otimes m}\).
Projection on the other parties decreases the unnormalized reduced
operator in positive-semidefinite order. More explicitly, partial
trace is cyclic with respect to operators on the traced subsystem,
and \(0\le P_{\ne i}\le I\). Therefore
\[
 q_m\sigma_i^{(m)}
 \le P_i^{(m)}\rho_i^{\otimes m}P_i^{(m)}
 \le 2^{-m(H_i-\delta)}P_i^{(m)}.
\]
Taking operator norms proves the assertion. This is a standard
typical-subspace construction, written out to make the simultaneous
marginal bound explicit.
\end{proof}

\begin{theorem}[Regularization identities]\label{thm:regular}
For every nonzero complex tripartite tensor \(T\),
\begin{align}
 \lim_{m\to\infty}F_{\theta,\infty}(T^{\otimes m})^{1/m}
 &=F_\theta(T),\label{eq:Freg}\\
 \lim_{m\to\infty}G_{\xi,\infty}(T^{\otimes m})^{1/m}
 &=G_\xi^{H}(T).\label{eq:Greg}
\end{align}
\end{theorem}
\begin{proof}
Choose any normalized representative \(\psi\) in the projective
orbit closure of \(T\). The outputs in Lemma \ref{lem:typical}
belong to the orbit closure of \(T^{\otimes m}\), since singular
local maps are limits of invertible ones. Applying the lemma,
then taking \(m\to\infty\), \(\delta\downarrow0\), and optimizing
over \(\psi\), gives the lower limits in both identities.
For \eqref{eq:Greg}, division by active \(\xi_i\) causes no problem:
their minimum is positive for the fixed vector \(\xi\).

For the first upper limit, use \(\Hinf\le H\) and multiplicativity:
\[
 F_{\theta,\infty}(T^{\otimes m})
 \le F_\theta(T^{\otimes m})=F_\theta(T)^m .
\]
For every \(\theta\in\Theta_3(\xi)\), a weighted average dominates
the minimum of the active ratios, hence
\[
 G_{\xi,\infty}(T^{\otimes m})
 \le F_{\theta,\infty}(T^{\otimes m})^{1/\langle\theta,\xi\rangle}
 \le F_\theta(T)^{m/\langle\theta,\xi\rangle}.
\]
Minimize over \(\theta\) and use the valid Shannon minimax identity.
The lower and upper limits agree.
\end{proof}

The unweighted identity \eqref{eq:Greg} already follows from
Derksen's constant-factor comparison of G-stable rank with slice
rank \cite[Proposition 4.9 and Theorem 5.2]{Derksen}, together with the known asymptotic slice-rank
formula \cite{CLZ}. The proof above also treats nonuniform weights
and \eqref{eq:Freg}; whether these formulations have appeared
elsewhere requires further priority checking.

A concrete three-copy restriction illustrates why the single-copy
value cannot be treated as multiplicative. In copy \(j\) of \(W_3\),
project party \(j\) onto \(\ket0\), for \(j=1,2,3\).
The result is a triangle of three entangled pairs. Each grouped
party has two maximally mixed qubits and one pure qubit, so
\[
 G_{\one,\infty}(W_3^{\otimes3})\ge4>(3/2)^3.
\]
This familiar entangled-pair construction is an illustration,
not a claimed new protocol. The three relevant scales are
\begin{center}
\begin{tabular}{lc}
\toprule
Operation on \(W_3\) & Value\\
\midrule
One-copy min-entropy radial value & \(3/2\)\\
Convexified one-copy entropy image & \(2^{2/3}\)\\
Tensor-regularized radial value & \(2^{H(1/3,\,2/3)}=3/2^{2/3}\)\\
\bottomrule
\end{tabular}
\end{center}
For completeness, the last value follows from Lemma \ref{lem:orbit}:
some smaller marginal eigenvalue is at most \(1/3\), and binary
entropy increases on \([0,1/2]\). The symmetric W state attains
the resulting upper bound. This is the established W asymptotic
slice-rank value, recovered here for comparison.

\section{Verification and limits of the claim}
The proofs above are mathematical arguments; finite computation
is supplementary. The accompanying deterministic script,
with random seed \(2026101707\), checks:
\begin{itemize}
 \item the rational contradiction \(27/8<4\) and the exact three-copy
 restriction, whose unnormalized marginal matrices have four
 diagonal entries equal to two and four equal to zero;
 \item the orbit inequality and weighted functional bound on
 \(200\) random local transforms for \(3\le k\le7\);
 \item the weighted scalar-root formula against independent linear
 feasibility calculations and explicit attaining states in \(32\)
 cases for \(3\le k\le6\);
 \item \(24\) finite instances of the positive-semidefinite domination
 in Lemma \ref{lem:typical}.
\end{itemize}
All checks passed. The largest root-versus-feasibility discrepancy
was \(1.83\times10^{-9}\) or less. Random tests do not prove a universal
statement, and these are author-context diagnostics, not independent
peer review or formal proof-assistant verification.

The additional R\'enyi verifier checks the order-three polynomial
gap and eight marginal matrices exactly, four symbolic curvature
identities, 48 concavity instances, and nine midpoint violations
above order two. The case \(\alpha=2.01\) uses \(a=10^{-400}\)
and 890 decimal digits, making explicit the precision needed near
the threshold. These finite checks supplement the proof of
Theorem \ref{thm:criterion}; they do not establish its universal
quantifiers by sampling.

The minimality verifier checks six symbolic or rational identities,
including the order-three curvature formula and the order-four
three-qubit witness. With seed \(2026101710\), it checks
\(100000\) boundary curvature cases and \(50000\) Jensen inequalities
for the complete boundary graph. The smallest numerical Jensen
gap is zero. The global convexity argument in
Theorem~\ref{thm:minimality}, including its endpoint treatment,
is the proof; finite tests are supplementary diagnostics.
The source archive accompanying this exact version on its
\href{https://airr.science/papers/ARR-2026-68GTBBQT1T83QATJ/versions/v1/}{AIRR record}
contains all three current verifiers, recorded outputs, dependency
versions, file hashes and the preserved earlier-version assessments.

The arXiv PDF's definitions and Theorem 60 were visually checked
against the source. A separate model review also downloaded the journal repository PDF
and visually checked printed page 327, confirming the identity and
its preceding concavity premise. The correction
claim concerns these two specified versions; a separate corrigendum
has not been ruled out. Neither the known value of the W
G-stable rank, the W-polytope geometry, classical convex duality,
nor the unweighted asymptotic identification is claimed as new.
The classical curvature of binary R\'enyi entropy is also credited;
the proposed new structural statements are the four-qubit criterion,
its minimality even within the R\'enyi family, and its sharp minimax
consequence. Their priority remains unchecked
beyond the targeted sources and searches described in the audit.
No result here concerns the general validity of the main Shannon
theory. An independent mathematical assessment and a broader
priority search are required before assigning significance or
asserting that this obstruction has not been noticed before.

\paragraph{AI assistance and review status.}
The drafting, literature inspection, calculations and computational
diagnostics were developed with substantial OpenAI Codex assistance.
Earlier nine- and eleven-page versions received model reviews
requested from GPT-6 Astra at xhigh effort. Those reports are
preserved with their exact PDF hashes. The order-three minimality
argument was subsequently examined as a separate supplement.
Any assessment of the final manuscript is identified by its exact
PDF hash in the accompanying record. Reviewer context is reused;
the model family has participated in the wider programme, and
statistical independence, certified context isolation and independently
observed runtime identity are not claimed. No human peer review,
formal proof-assistant verification or editorial acceptance is implied.

\begin{thebibliography}{9}
\bibitem{CLZ}
M. Christandl, V. Lysikov and J. Zuiddam,
\emph{Weighted slice rank and a minimax correspondence to Strassen's spectra},
J. Math. Pures Appl. \textbf{172} (2023), 299--329.
Versions assessed: \href{https://arxiv.org/abs/2012.14412v3}{arXiv:2012.14412v3}
(19 April 2023), Theorem 60 on p.~26; and the
\href{https://dare.uva.nl/id/5f9e575f-cd5d-4316-8f0c-9320c049778f}{final journal version},
Theorem 60 on p.~327.
\href{https://doi.org/10.1016/j.matpur.2023.02.006}{doi:10.1016/j.matpur.2023.02.006}.

\bibitem{Derksen}
H. Derksen, \emph{The G-stable rank for tensors and the cap set problem},
Algebra Number Theory \textbf{16} (2022), 1071--1097.
\href{https://arxiv.org/abs/2002.08435}{arXiv:2002.08435}.
\href{https://doi.org/10.2140/ant.2022.16.1071}{doi:10.2140/ant.2022.16.1071}.

\bibitem{CVZ}
M. Christandl, P. Vrana and J. Zuiddam,
\emph{Universal points in the asymptotic spectrum of tensors},
J. Amer. Math. Soc. \textbf{36} (2023), 31--79.
\href{https://arxiv.org/abs/1709.07851}{arXiv:1709.07851}.
\href{https://doi.org/10.1090/jams/996}{doi:10.1090/jams/996}.

\bibitem{Walter}
M. Walter, B. Doran, D. Gross and M. Christandl,
\emph{Entanglement polytopes: Multiparticle entanglement from single-particle information},
Science \textbf{340} (2013), 1205--1208.
\href{https://arxiv.org/abs/1208.0365}{arXiv:1208.0365}.
\href{https://doi.org/10.1126/science.1232957}{doi:10.1126/science.1232957}.
\bibitem{DVC}
W. D\"ur, G. Vidal and J. I. Cirac,
\emph{Three qubits can be entangled in two inequivalent ways},
Phys. Rev. A \textbf{62} (2000), 062314, Sections II--III.
\href{https://arxiv.org/abs/quant-ph/0005115v2}{arXiv:quant-ph/0005115v2}.
\href{https://doi.org/10.1103/PhysRevA.62.062314}{doi:10.1103/PhysRevA.62.062314}.

\bibitem{Hillion}
E. Hillion and O. Johnson,
\emph{A proof of the Shepp--Olkin entropy concavity conjecture},
Bernoulli \textbf{23} (2017), 3638--3649.
\href{https://arxiv.org/abs/1503.01570}{arXiv:1503.01570}.
\href{https://doi.org/10.3150/16-BEJ860}{doi:10.3150/16-BEJ860}.
\end{thebibliography}
\end{document}
